\documentclass[12pt]{article}
\usepackage[T1]{fontenc}
\usepackage[utf8]{inputenc}
\usepackage{lmodern}
\usepackage{textcomp}
\usepackage{enumitem}
\usepackage{geometry}
\usepackage{graphicx}
\usepackage{amsmath, amssymb}
\geometry{margin=1in}
\begin{document}
\begin{center}
Sociology - Graduate Level Assessment 17 \
Total Marks: 40 \
Answer all questions.
\end{center}

\section*{Multiple Choice (10 marks)}
\begin{enumerate}[label=\arabic*.]
\item The term 'collective consumption' (Castells 1977) refers to:
\begin{enumerate}[label=(\alph*)]
    \item the privatization of public services by the Conservative government
    \item the lifestyle practice of shopping in peer groups
    \item the form of tuberculosis suffered by those who collect stamps
    \item the provision of health, housing, and education services by the state
\end{enumerate}
\item Religious organizations such as the Church of Norway, Islam, the Church of England, and the Church of Greece have which of the following characteristics in common?
\begin{enumerate}[label=(\alph*)]
    \item They are cultures as well as churches.
    \item They practice separation of church and state.
    \item They exclude women as clergy.
    \item They are monotheistic.
\end{enumerate}
\item Scott (1991) introduced the term 'power elite' to describe:
\begin{enumerate}[label=(\alph*)]
    \item the ruling class, or bourgeoisie, who exploit the proletariat
    \item a capitalist class dependent on property ownership and advantaged life chances
    \item an alignment of classes with shared interests but no state power
    \item a state elite whose members are drawn overwhelmingly from a power bloc
\end{enumerate}
\item Warner and Myrdal claimed that black former slaves were not joining the 'ethnic melting pot' of the North American cities because:
\begin{enumerate}[label=(\alph*)]
    \item the white population did not believe in liberty, equality, and democracy
    \item they wanted to retain a strong sense of their original ethnic identity
    \item they were not prepared to leave the Southern states and move to North America
    \item the promise of citizenship was contradicted by continued discrimination
\end{enumerate}
\item An ecclesia is:
\begin{enumerate}[label=(\alph*)]
    \item a religious organization that claims total spiritual authority over its members
    \item a church organized around voluntary rather than compulsory membership
    \item a sect or cult with a very small following
    \item a hierarchy of priests or other spiritual leaders
\end{enumerate}
\end{enumerate}

\section*{True / False (10 marks)}
\textit{Answer True or False}
\begin{enumerate}[label=\arabic*.]
\setcounter{enumi}{5}
\item True or False: The office is not recognized as a level of society, unlike the household, nation state, and global village.
\item True or False: Mead's 'generalized other' is defined as a significant figure in childhood who imparts the general values of society.
\item True or False: Booth's (1901) study suggested that 27.50\% of Londoners were experiencing poverty at that time.
\item True or False: Marx's concept of class polarization refers to the increasing gap between the rich and poor, leading to heightened class consciousness.
\item True or False: War became possible between nation states in the nineteenth century due to rival empires, industrialization, and intense competition.
\end{enumerate}

\section*{Long Form Response (20 marks)}
\textit{Answer each question with a clear and complete explanation.}
\begin{enumerate}[label=\arabic*.]
\setcounter{enumi}{10}
\item Discuss the role of informal negative sanctions in shaping social behavior, using the example of children who are teased for thumb-sucking upon entering kindergarten.
\item Discuss the naturalistic explanations of race proposed by nineteenth-century scientific theories, and critically analyze which aspects were commonly emphasized or overlooked in these discussions.
\end{enumerate}

\vfill
\noindent\textit{End of Paper}
\end{document}